% ===========================================================================
%  Worksheet --- Automorphisms, Irreducibility and Splitting Fields
%  Problems 7--12: a tour through the roots of unity.
%
%  THE WRITE-UP, and it is the assistant's to type.
%    The assistant transcribes each question or topic faithfully and emits an
%    empty, marked solution region beneath it. The student does the
%    mathematics; the assistant typesets what they agreed was right.
%
%  ONE `problem` ENVIRONMENT PER LETTERED PART, each with its own matched
%  SOLUTION / END SOLUTION pair. `board hw status` reads those markers to say
%  what is written up and what is next, and it pairs them by label.
%
%  Problems 7, 8, 9 and 10 are the four chosen for this sitting. Problems 11
%  and 12 are the finite-field pair and are not stated here; they get their own
%  sheet when they are worked.
% ===========================================================================
\documentclass[11pt]{article}
\usepackage{coursemacros}

\title{Worksheet --- Automorphisms, Irreducibility and Splitting Fields}
\author{Mikey Ferguson}
\date{\today}

\begin{document}
\maketitle

\noindent
Throughout, $\zeta_n = e^{2\pi i/n}$, and $\operatorname{Aut}(K/F)$ denotes the
group of automorphisms of $K$ that fix $F$ pointwise.

% One problem/solution pair per thing worked, in the sheet's own order. A
% region stays empty until the answer in it is agreed correct.

\begin{problem}{7(a)}
  \emph{Automorphisms move roots to roots.} Let $K/F$ be an extension and let
  $\sigma \in \operatorname{Aut}(K/F)$. Show that if $f \in F[x]$ and
  $f(\alpha) = 0$ for some $\alpha \in K$, then $f(\sigma(\alpha)) = 0$.
\end{problem}

% ===== SOLUTION 7(a) =====
\begin{proof}
Write $f = \sum_j c_j x^j$ with every $c_j \in F$, so that
$f(\alpha) = \sum_j c_j \alpha^j = 0$. Then
\begin{align*}
  f(\sigma(\alpha))
    &= \sum_j c_j\, \sigma(\alpha)^j \\
    &= \sum_j c_j\, \sigma(\alpha^j) && \text{($\sigma$ multiplies)} \\
    &= \sum_j \sigma(c_j \alpha^j) && \text{($\sigma$ fixes $c_j \in F$, and multiplies)} \\
    &= \sigma\Bigl(\sum_j c_j \alpha^j\Bigr) && \text{($\sigma$ adds)} \\
    &= \sigma(0) = 0 && \text{(since $0 \in F$)}. \qedhere
\end{align*}
\end{proof}
% ===== END SOLUTION 7(a) =====

\begin{problem}{7(b)}
  Let $K/F$ be an extension and let $\sigma \in \operatorname{Aut}(K/F)$. Show
  that if $K = F(\alpha_1, \dots, \alpha_r)$, then $\sigma$ is determined by
  $\sigma(\alpha_1), \dots, \sigma(\alpha_r)$.
\end{problem}

% ===== SOLUTION 7(b) =====
\begin{proof}
Let $\tau \in \operatorname{Aut}(K/F)$ with $\tau(\alpha_i) = \sigma(\alpha_i)$
for every $i$. We prove $\tau = \sigma$. Write $\alpha$ for
$(\alpha_1, \dots, \alpha_r)$.

\emph{Polynomials.} If $p \in F[x_1, \dots, x_r]$, then $p(\alpha)$ is a sum of
terms $c\,\alpha_1^{e_1}\cdots\alpha_r^{e_r}$ with $c \in F$. Since $\sigma$ and
$\tau$ add, multiply, fix $c$, and agree on every $\alpha_i$, we get
$\sigma(p(\alpha)) = \tau(p(\alpha))$.

\emph{Quotients.} Let $x \in K$. Then $x = p(\alpha)\,q(\alpha)^{-1}$ with
$p, q \in F[x_1, \dots, x_r]$ and $q(\alpha) \neq 0$, so the inverse exists in
$K$. Then
\begin{align*}
  \sigma(x)
    &= \sigma\bigl(p(\alpha)\,q(\alpha)^{-1}\bigr) \\
    &= \sigma(p(\alpha))\,\sigma\bigl(q(\alpha)^{-1}\bigr) && \text{($\sigma$ multiplies)} \\
    &= \sigma(p(\alpha))\,\sigma(q(\alpha))^{-1} && \text{($\sigma$ preserves inverses)} \\
    &= \tau(p(\alpha))\,\tau(q(\alpha))^{-1} && \text{(polynomials)} \\
    &= \tau(p(\alpha))\,\tau\bigl(q(\alpha)^{-1}\bigr) && \text{($\tau$ preserves inverses)} \\
    &= \tau\bigl(p(\alpha)\,q(\alpha)^{-1}\bigr) = \tau(x) && \text{($\tau$ multiplies)}.
\end{align*}
So $\sigma = \tau$, and $\sigma$ is determined by
$\sigma(\alpha_1), \dots, \sigma(\alpha_r)$.
\end{proof}
% ===== END SOLUTION 7(b) =====

\begin{problem}{7(c)}
  Let $K = \mathbb{Q}(\zeta_n)$. Show that every
  $\sigma \in \operatorname{Aut}(K/\mathbb{Q})$ satisfies
  $\sigma(\zeta_n) = \zeta_n^{\,k}$ for some $k$ with $\gcd(k,n) = 1$, and that
  $\sigma \mapsto k \bmod n$ is an injective group homomorphism
  \[
    \operatorname{Aut}(K/\mathbb{Q}) \hookrightarrow (\mathbb{Z}/n\mathbb{Z})^{\times} .
  \]
  In particular $\lvert \operatorname{Aut}(K/\mathbb{Q}) \rvert \leq \varphi(n)$.
\end{problem}

% ===== SOLUTION 7(c) =====
\begin{proof}[Proof that $\gcd(k,n)=1$]
Let $\sigma(\zeta_n)=\zeta_n^{\,k}$, and suppose instead that
$d=\gcd(k,n)>1$. Then $n/d$ and $k/d$ are integers, and
\[
  \zeta_n^{\,n/d} = e^{\frac{2\pi i}{n}\cdot\frac{n}{d}} = e^{2\pi i/d} \neq 1,
\]
since $d>1$ puts $2\pi/d$ in $(0,2\pi)$. On the other hand,
\begin{align*}
  \sigma\bigl(\zeta_n^{\,n/d}\bigr)
    &= \sigma(\zeta_n)^{n/d}
     = \bigl(\zeta_n^{\,k}\bigr)^{n/d}
     = \bigl(\zeta_n^{\,n}\bigr)^{k/d}
     = 1 = \sigma(1).
\end{align*}
Since $\sigma$ is injective, $\zeta_n^{\,n/d}=1$, a contradiction. So
$\gcd(k,n)=1$.
\end{proof}
\begin{proof}[Proof that $\sigma\mapsto k \bmod n$ is a homomorphism]
Let $\sigma(\zeta_n)=\zeta_n^{\,k}$ and $\tau(\zeta_n)=\zeta_n^{\,l}$. Then
\begin{align*}
  (\sigma\circ\tau)(\zeta_n)
    &= \sigma\bigl(\tau(\zeta_n)\bigr)
     = \sigma\bigl(\zeta_n^{\,l}\bigr)
     = \sigma(\zeta_n)^{l}
     = \bigl(\zeta_n^{\,k}\bigr)^{l}
     = \zeta_n^{\,kl} \\
    &= e^{2\pi i\,kl/n}
     = e^{2\pi i\,(kl \bmod n)/n}
     = \zeta_n^{\,kl \bmod n},
\end{align*}
where $\sigma(\zeta_n^{\,l})=\sigma(\zeta_n)^{l}$ is
$\sigma(ab)=\sigma(a)\sigma(b)$ applied $l-1$ times. So
$\sigma\circ\tau\mapsto kl \bmod n$, and composition goes to multiplication
mod $n$.
\end{proof}
\begin{proof}[Proof that $\sigma\mapsto k \bmod n$ is injective]
Let $\sigma(\zeta_n)=\zeta_n^{\,k}$ and $\tau(\zeta_n)=\zeta_n^{\,k'}$, and
suppose $k\equiv k' \pmod n$, so $k'=k+hn$ for some $h\in\mathbb{Z}$. Then
\[
  \zeta_n^{\,k'} = e^{2\pi i(k+hn)/n}
    = e^{2\pi i k/n}\, e^{2\pi i h}
    = e^{2\pi i k/n} = \zeta_n^{\,k},
\]
so $\sigma(\zeta_n)=\tau(\zeta_n)$. Since $\zeta_n$ is algebraic over
$\mathbb{Q}$, $K=\mathbb{Q}(\zeta_n)=\mathbb{Q}[\zeta_n]$, so every $x\in K$
has the form $x=\sum_{j=0}^{m} a_j\zeta_n^{\,j}$ with $a_j\in\mathbb{Q}$.
Because $\sigma$ and $\tau$ are automorphisms of $K$ (additive and
multiplicative) fixing each $a_j\in\mathbb{Q}$,
\[
  \sigma(x) = \sum_{j=0}^{m} a_j\,\sigma(\zeta_n)^{j}
    = \sum_{j=0}^{m} a_j\bigl(\zeta_n^{\,k}\bigr)^{j}
    = \sum_{j=0}^{m} a_j\bigl(\zeta_n^{\,k'}\bigr)^{j}
    = \sum_{j=0}^{m} a_j\,\tau(\zeta_n)^{j}
    = \tau(x).
\]
Hence $\sigma=\tau$, and the map is injective.
\end{proof}
\begin{proof}[Proof that $\lvert\mathrm{Aut}(K/\mathbb{Q})\rvert\le\varphi(n)$]
$\varphi(n)$ is the number of $k\in\{1,\dots,n\}$ with $\gcd(k,n)=1$.
Suppose $\lvert\mathrm{Aut}(K/\mathbb{Q})\rvert>\varphi(n)$. Every
automorphism sends $\zeta_n$ to $\zeta_n^{\,k}$ with $\gcd(k,n)=1$, so there
are $\sigma_1\ne\sigma_2$ with $\sigma_1(\zeta_n)=\zeta_n^{\,k}=\sigma_2(\zeta_n)$
for the same $k$. From earlier, though, $\sigma_1(\zeta_n)=\sigma_2(\zeta_n)$
forces $\sigma_1=\sigma_2$, a contradiction.
\end{proof}
% ===== END SOLUTION 7(c) =====

\begin{problem}{8(a)}
  \emph{Conjugate roots and the cyclotomic polynomial $\Phi_p$.} Let
  $p \in F[x]$ be irreducible, with roots $\alpha$ and $\beta$ in (possibly
  different) extensions of $F$. Show there is a unique isomorphism
  $F(\alpha) \to F(\beta)$ fixing $F$ and sending $\alpha \mapsto \beta$.
\end{problem}

% ===== SOLUTION 8(a) =====
\begin{proof}[Existence]
Let $H : F[x] \to F(\alpha)$ and $G : F[x] \to F(\beta)$ be the evaluation
maps $f \mapsto f(\alpha)$ and $f \mapsto f(\beta)$. Both are surjective,
since every element of $F(\alpha)$ is a polynomial in $\alpha$, and likewise
for $\beta$. Since $p(\alpha) = 0$, we have $(p) \subseteq \ker H$; since $p$ is
irreducible, $(p)$ is maximal, and $\ker H \neq F[x]$, so $\ker H = (p)$.
Likewise $\ker G = (p)$. By the first isomorphism theorem there are
isomorphisms
\[
  g : F(\alpha) \to F[x]/(p), \qquad h : F(\beta) \to F[x]/(p),
\]
with $g(f(\alpha)) = f + (p)$ and $h(f(\beta)) = f + (p)$. Put
$\sigma = h^{-1} \circ g : F(\alpha) \to F(\beta)$. It is an isomorphism, it
fixes $F$, and $\sigma(\alpha) = h^{-1}(x + (p)) = \beta$.
For example, over $\mathbb{Q}$ with $p = x^2 - 2$, this gives
$\sqrt{2} \mapsto -\sqrt{2}$.
\end{proof}
\begin{proof}[Uniqueness]
Let $\tau : F(\alpha) \to F(\beta)$ fix $F$ with $\tau(\alpha) = \beta$. For
$c \in F$, $\tau(c) = c = \sigma(c)$. Let $x \in F(\alpha)$, so
$x = \sum_{i=0}^{n} a_i \alpha^i$ with $a_i \in F$. Then
\begin{align*}
  \tau(x) &= \sum_{i=0}^{n} a_i\, \tau(\alpha)^i && \text{($\tau$ adds, multiplies, fixes $F$)} \\
          &= \sum_{i=0}^{n} a_i\, \beta^i \\
          &= \sum_{i=0}^{n} a_i\, \sigma(\alpha)^i = \sigma(x) && \text{(likewise for $\sigma$)}.
\end{align*}
So $\tau = \sigma$.
\end{proof}
% ===== END SOLUTION 8(a) =====

\begin{problem}{8(b)}
  Let $p$ be prime and
  $\Phi_p(x) = x^{p-1} + x^{p-2} + \cdots + x + 1$. Show that $\Phi_p$ is
  irreducible over $\mathbb{Q}$.
\end{problem}

% ===== SOLUTION 8(b) =====
\begin{proof}
We have $(x-1)\Phi_p(x) = x^p - 1$. Substituting $x+1$ for $x$,
\[
  x\,\Phi_p(x+1) = (x+1)^p - 1
    = \Bigl(\sum_{i=0}^{p} \binom{p}{i} x^i\Bigr) - 1
    = \sum_{i=1}^{p} \binom{p}{i} x^i ,
\]
so
\[
  \Phi_p(x+1) = \sum_{i=1}^{p} \binom{p}{i} x^{i-1}
    = x^{p-1} + \binom{p}{p-1} x^{p-2} + \cdots + \binom{p}{1}.
\]
This is irreducible by Eisenstein at $p$: $p \mid \binom{p}{1} = p$ and
$p^2 \nmid \binom{p}{1}$; $p \mid \binom{p}{k}$ for $k \in \{2, \dots, p-1\}$;
and $p \nmid \binom{p}{p} = 1$. If $\Phi_p(x)$ were reducible, so would
$\Phi_p(x+1)$ be. So $\Phi_p(x)$ is irreducible.
\end{proof}
% ===== END SOLUTION 8(b) =====

\begin{problem}{8(c)}
  Deduce that $[\mathbb{Q}(\zeta_p) : \mathbb{Q}] = p - 1$ and that the
  injection of Problem 7(c) is an isomorphism:
  \[
    \operatorname{Aut}(\mathbb{Q}(\zeta_p)/\mathbb{Q})
      \cong (\mathbb{Z}/p\mathbb{Z})^{\times} .
  \]
\end{problem}

% ===== SOLUTION 8(c) =====
\begin{proof}[Proof that {$[\mathbb{Q}(\zeta_p):\mathbb{Q}] = p-1$}]
Let $\zeta_p = e^{2\pi i/p}$. We have $x^p - 1 = (x-1)\Phi_p(x)$, and
$\zeta_p$ is a root of $x^p - 1$ but $\zeta_p \neq 1$, so $\zeta_p$ is a
root of $\Phi_p$. By 8(b), $\Phi_p$ is irreducible over $\mathbb{Q}$, and it
is monic, so it is the minimal polynomial of $\zeta_p$. Hence
$[\mathbb{Q}(\zeta_p):\mathbb{Q}] = \deg \Phi_p = p-1$.
\end{proof}

% TODO(mferguson): that the injection of 7(c) is onto is not yet proved.
% ===== END SOLUTION 8(c) =====

\begin{problem}{9(a)}
  \emph{The eighth roots of unity and $x^4 + 1$.} Let $\zeta = \zeta_8$ and
  $K = \mathbb{Q}(\zeta)$. Show that $\zeta$ is a root of $x^4 + 1$, that
  $x^4 + 1$ is irreducible over $\mathbb{Q}$, and hence that
  $[K : \mathbb{Q}] = 4$.
\end{problem}

% ===== SOLUTION 9(a) =====
% TODO(mferguson): your work goes here.
% ===== END SOLUTION 9(a) =====

\begin{problem}{9(b)}
  With $\zeta = \zeta_8$ and $K = \mathbb{Q}(\zeta)$, show that $x^4 + 1$
  factors into quadratics over each of $\mathbb{Q}(\sqrt{2})$, $\mathbb{Q}(i)$
  and $\mathbb{Q}(\sqrt{-2})$.
\end{problem}

% ===== SOLUTION 9(b) =====
% TODO(mferguson): your work goes here.
% ===== END SOLUTION 9(b) =====

\begin{problem}{9(c)}
  With $\zeta = \zeta_8$ and $K = \mathbb{Q}(\zeta)$, show that
  $K = \mathbb{Q}(i, \sqrt{2})$.
\end{problem}

% ===== SOLUTION 9(c) =====
% TODO(mferguson): your work goes here.
% ===== END SOLUTION 9(c) =====

\begin{problem}{9(d)}
  With $\zeta = \zeta_8$ and $K = \mathbb{Q}(\zeta)$: for $k \in \{3,5,7\}$,
  show there is an automorphism $\sigma_k$ of $K$ with
  $\sigma_k(\zeta) = \zeta^{k}$, that each $\sigma_k$ has order $2$, and that
  the subfields of $K$ fixed by $\sigma_5$, $\sigma_7$ and $\sigma_3$ are
  $\mathbb{Q}(i)$, $\mathbb{Q}(\sqrt{2})$ and $\mathbb{Q}(\sqrt{-2})$
  respectively.
\end{problem}

% ===== SOLUTION 9(d) =====
% TODO(mferguson): your work goes here.
% ===== END SOLUTION 9(d) =====

\begin{problem}{10(a)}
  \emph{A first splitting field: $x^3 - 2$.} Let $\omega = \zeta_3$. Show that
  the complex roots of $x^3 - 2$ are $\sqrt[3]{2}$, $\omega\sqrt[3]{2}$ and
  $\omega^2\sqrt[3]{2}$, and conclude that its splitting field over
  $\mathbb{Q}$ is $E = \mathbb{Q}(\sqrt[3]{2}, \omega)$.
\end{problem}

% ===== SOLUTION 10(a) =====
% TODO(mferguson): your work goes here.
% ===== END SOLUTION 10(a) =====

\begin{problem}{10(b)}
  With $\omega = \zeta_3$ and $E = \mathbb{Q}(\sqrt[3]{2}, \omega)$, show that
  $[E : \mathbb{Q}] = 6$. In particular, $\mathbb{Q}(\sqrt[3]{2})$ contains a
  root of $x^3 - 2$ but is not its splitting field.
\end{problem}

% ===== SOLUTION 10(b) =====
% TODO(mferguson): your work goes here.
% ===== END SOLUTION 10(b) =====

\begin{problem}{10(c)}
  More generally, let $a \in \mathbb{Q}^{\times}$ and let $\alpha$ be any
  complex root of $x^n - a$. Show that the splitting field of $x^n - a$ over
  $\mathbb{Q}$ is $\mathbb{Q}(\alpha, \zeta_n)$.
\end{problem}

% ===== SOLUTION 10(c) =====
% TODO(mferguson): your work goes here.
% ===== END SOLUTION 10(c) =====

\begin{problem}{10(d)}
  With $E = \mathbb{Q}(\sqrt[3]{2}, \zeta_3)$, show that every
  $\sigma \in \operatorname{Aut}(E/\mathbb{Q})$ permutes the three roots of
  $x^3 - 2$ and is determined by that permutation. Conclude that
  $\operatorname{Aut}(E/\mathbb{Q})$ is isomorphic to a subgroup of $S_3$.
\end{problem}

% ===== SOLUTION 10(d) =====
% TODO(mferguson): your work goes here.
% ===== END SOLUTION 10(d) =====

\end{document}
